\section{Backend dispatch}

\begin{frame}{One dispatcher, two CLI backends}
  \begin{center}
  \begin{tikzpicture}[node distance=0.7cm and 1.0cm,
                      every node/.style={font=\footnotesize}]
    \node[agent3, minimum width=3.4cm] (prop)
      {\cmd{\_propose\_params}\\\scriptsize \cmd{digital\_autoresearch.py:504}};
    \node[notebox, above=0.4cm of prop, minimum width=3.4cm]
      (sw) {switch \cmd{self.backend}};
    \node[bxclaude,  below left =1.0cm and 0.4cm of prop] (cc)
      {Claude Code CLI\\\scriptsize Opus 4.7, subscription};
    \node[bxopencode, below right=1.0cm and 0.4cm of prop] (oc)
      {opencode CLI\\\scriptsize \cmd{openai/gpt-5.3-codex}};
    \node[bxclaude,  below=of cc, minimum width=3.4cm]
      (cch) {\cmd{ClaudeCodeHarness}\\\scriptsize \cmd{claude\_code\_harness.py:207}};
    \node[bxopencode, below=of oc, minimum width=3.4cm]
      (och) {\cmd{OpenCodeHarness}\\\scriptsize \cmd{opencode\_harness.py:123}};
    \draw[arrow] (sw) -- (prop);
    \draw[arrowclaude] (prop) -- (cc);
    \draw[arrowopen]   (prop) -- (oc);
    \draw[arrowclaude] (cc)  -- (cch);
    \draw[arrowopen]   (oc)  -- (och);
  \end{tikzpicture}
  \end{center}
  \vspace{0.2em}
  \footnotesize
  \begin{itemize}
    \item Same async \cmd{run() $\rightarrow$ HarnessResult}
          contract on both sides; the loop never branches.
    \item Same dispatch shape as \cmd{idea\_to\_rtl.py:\_build\_harness:326}
          (the original pattern). LiteLLM is a dev-only fallback;
          not part of the demo.
    \item OAuth handled by each CLI; no provider keys in the
          subprocess environment.
  \end{itemize}
  \note{%
    Why this matters: the loop body is unchanged across backends,
    so whatever bugs exist in the loop are independent of which CLI
    proposes the next params. The harness contract is also what
    lets us do an A/B run: same RTL, same testbench, same FoM,
    same budget; only the \emph{policy} (the LLM behind the CLI)
    changes. Each CLI is invoked via \cmd{run --print} or
    \cmd{run --format json} so the agent inherits the user's
    subscription session, not an API key.}
\end{frame}

\begin{frame}[fragile]{Why dispatch matters: the bug \cmd{f1b6d13} fixed}
  \begin{lstlisting}[style=pystyle]
# digital_autoresearch.py:_propose_params  (post-fix)
if self.backend == "cc_cli":
    return await self._propose_params_via_cc_cli(...)
if self.backend == "opencode":
    return await self._propose_params_via_opencode(...)
# legacy fallback (development only)
return await litellm.acompletion(...)
  \end{lstlisting}
  \footnotesize
  \begin{itemize}
    \item \warn{Pre-fix bug:} \cmd{--backend opencode --model
          openai/gpt-5.3-codex} silently routed flow-strategy
          proposals through LiteLLM straight to OpenAI, bypassing
          the user's opencode CLI subscription.
    \item Same bug shape was latent for \cmd{cc\_cli}: flow-strategy
          proposals defaulted to whatever LiteLLM model was set, not
          Opus under the user's Claude subscription.
    \item Regression guard:
          \cmd{tests/test\_digital\_autoresearch.py::TestProposalDispatch:881}
          asserts that each backend instantiates the correct harness
          class.
  \end{itemize}
  \note{%
    The fix is three lines, but the failure mode was invisible
    without staring at network traffic. Flow-strategy proposals
    pre-dated cc\_cli/opencode support and were never updated to
    dispatch by \cmd{self.backend}; only \cmd{rtl} and \cmd{hybrid}
    strategies were routed correctly. The lesson is that
    cross-cutting changes (a new backend) must touch every place
    that talks to the model, not just the most obvious one.
    \cmd{TestProposalDispatch} now nails the contract so this
    cannot regress without test failure.}
\end{frame}
